\usepackage[english]{babel}
\usepackage{kotex}    

\usepackage[letterpaper,top=1.2cm,bottom=1.2cm,left=1.2cm,right=1.2cm,marginparwidth=1.75cm]{geometry}

% Useful packages
\usepackage{amsmath, amsthm, amssymb, graphicx, enumitem, authblk, tikz, tikz-cd, verbatim, relsize}
\usepackage{thmtools}
\usepackage{tcolorbox}
\usepackage{forest} % 트리 구조를 위한 패키지
\usepackage{mathtools}
\usepackage[colorlinks=true, allcolors=black]{hyperref}
\usepackage[framemethod=TikZ]{mdframed}
\usepackage{lmodern}

\DeclareMathAlphabet{\mathtt}{OT1}{lmtt}{m}{n}

\usepackage{fontspec}
\setmainfont{neodgm_code.ttf}[
  Path = ./,
  UprightFont = *,
  Ligatures = TeX,
  Scale = MatchLowercase,
  % 폰트에 Bold/Italic이 없을 때 임시로:
  AutoFakeBold = 2.0,
  AutoFakeSlant = 0.2
]

\usetikzlibrary{calc}
\usetikzlibrary{shapes,backgrounds}
\usetikzlibrary{decorations.pathreplacing}
\usetikzlibrary{arrows.meta, positioning, matrix, fit}

% \usetikzlibrary{arrows.meta, positioning, external}
% \pgfplotsset{compat=1.18} % Set compatibility for pgfplots

\newtheoremstyle{romanstyle} % 스타일 이름
  {10pt} % 위쪽 여백
  {10pt} % 아래쪽 여백
  {\rmfamily} % 본문 로마체
  {} % 들여쓰기 없음
  {\bfseries} % 제목 굵게
  {.} % 제목 뒤 구두점
  { } % 제목과 본문 사이의 기본 간격
  {} % 제목 레이아웃 기본값

% 새 스타일 적용
\theoremstyle{romanstyle}


% 1) 박스형 스타일 프리셋
\declaretheoremstyle[
  headfont=\bfseries, bodyfont=\rmfamily,
  headpunct={.}, postheadspace=.5em,
  mdframed={linewidth=0.6pt, linecolor=black, backgroundcolor=white,
            innerleftmargin=8pt, innerrightmargin=8pt,
            innertopmargin=6pt, innerbottommargin=6pt,
            skipabove=10pt, skipbelow=10pt, roundcorner=0pt}
]{thm-plain}

\declaretheoremstyle[
  headfont=\bfseries, bodyfont=\rmfamily, headpunct={.},
  postheadspace=.5em,
  mdframed={linewidth=1pt, linecolor=black, backgroundcolor=gray!4,
            leftline=true, rightline=false, topline=false, bottomline=false,
            innerleftmargin=10pt, innerrightmargin=8pt,
            innertopmargin=6pt, innerbottommargin=6pt,
            skipabove=10pt, skipbelow=10pt}
]{thm-left}

\declaretheoremstyle[
  headfont=\bfseries, bodyfont=\rmfamily, headpunct={.},
  postheadspace=.5em,
  mdframed={linewidth=0.8pt, linecolor=black, backgroundcolor=white,
            roundcorner=2pt, tikzsetting={dash pattern=on 3pt off 2pt},
            innerleftmargin=8pt, innerrightmargin=8pt,
            innertopmargin=6pt, innerbottommargin=6pt}
]{thm-dashed}

\declaretheoremstyle[
  headfont=\bfseries, bodyfont=\rmfamily, headpunct={.},
  postheadspace=.5em,
  mdframed={linewidth=0.8pt, linecolor=blue!60!black, backgroundcolor=black!3,
            roundcorner=1.5pt,
            innerleftmargin=8pt, innerrightmargin=8pt,
            innertopmargin=6pt, innerbottommargin=6pt}
]{thm-blue}

\declaretheorem[style=thm-left,   numberwithin=subsection]{theorem}
\declaretheorem[style=thm-blue,   numberwithin=subsection]{proposition}
\declaretheorem[style=thm-dashed, numberwithin=subsection]{lemma}
\declaretheorem[style=thm-blue,   numberwithin=subsection]{corollary}
\declaretheorem[style=thm-plain,  numberwithin=subsection]{definition}

\newtheorem{axiom}{Axiom}

% Unnumbered environments
\newtheorem*{observation}{Observation}
\newtheorem*{example}{Example}
\newtheorem*{solution}{Solution}

% Mathematical commands
\newcommand{\st}{ \ \operatorname{s.t} \ }
\newcommand{\as}{ \ \operatorname{as} \ }
\newcommand{\setx}{\{x_0, x_1, \dots, x_n\}}
\newcommand{\setn}{\{1, \dots, n\}}
\newcommand{\dis}{\displaystyle}
\newcommand{\lete}{Let $\epsilon > 0$ be given.}
\newcommand{\fs}{\text{ for some }}
\newcommand{\defequal}{\overset{\textnormal{def}}{=}}
\newcommand{\setequal}{\overset{\textnormal{set}}{=}}
\newcommand*{\aand}{_{\text{ and }}}
\newcommand*{\oor}{_{\text{ or }}}
\DeclareMathOperator{\argmax}{argmax}
\DeclareMathOperator{\argmin}{argmin}
\newcommand{\im}{\operatorname{im}}
\newcommand{\nullity}{\operatorname{nullity}}
\newcommand{\rank}{\operatorname{rank}}
\newcommand{\ipvs}{\langle \, \cdot \, \rangle}
\newcommand{\diam}{\operatorname{diam}} % Add this line in the preamble
\newlength{\datedwd}
\newcommand{\dated}[2]{%
  \settowidth{\datedwd}{#1}%
  \noindent
  \parbox[t]{\datedwd}{\strut #1}%
  \hspace{0.5em}%
  \begin{minipage}[t]{\dimexpr\linewidth-\datedwd-0.5em\relax}
    \vspace{-0.65em}%
    \begin{enumerate}[label=\arabic*., leftmargin=*, itemsep=0pt, topsep=0pt, parsep=0pt]
      #2
    \end{enumerate}
  \end{minipage}\par
}

% \addparttoc{제목} 형태로 사용하도록 정의
\newcommand{\addparttoc}[1]{%
    \phantomsection
    \addtocontents{toc}{\protect\addvspace{1.0\baselineskip}} % 목차 내 여백 추가
    \addcontentsline{toc}{part}{#1}% 목차에 'part' 레벨로 추가
}


% ===================== TOC에 정리 제목(옵션 있을 때만) 자동 추가 ===================== %
\usepackage{xparse}   %% <<< ADDED
\usepackage{etoolbox} %% <<< ADDED
\setcounter{tocdepth}{3}    %% <<< ADDED: subsubsection까지 TOC 표시
\setcounter{secnumdepth}{3} %% <<< ADDED: (선택) 번호도 3단계까지

% theorem
\let\theoremOLD\theorem           %% <<< ADDED
\let\endtheoremOLD\endtheorem     %% <<< ADDED
\RenewDocumentEnvironment{theorem}{o}{%% <<< ADDED
  \IfNoValueTF{#1}{%
    \theoremOLD
  }{%
    \theoremOLD
    \phantomsection
    \addcontentsline{toc}{subsubsection}{%
      \protect\numberline{\thetheorem}\ #1}%
  }%
}{\endtheoremOLD}

% lemma
\let\lemmaOLD\lemma
\let\endlemmaOLD\endlemma
\RenewDocumentEnvironment{lemma}{o}{%
  \IfNoValueTF{#1}{%
    \lemmaOLD
  }{%
    \lemmaOLD
    \phantomsection
    \addcontentsline{toc}{subsubsection}{%
      \protect\numberline{\thelemma}\ #1}%
  }%
}{\endlemmaOLD}

% proposition
\let\propositionOLD\proposition
\let\endpropositionOLD\endproposition
\RenewDocumentEnvironment{proposition}{o}{%
  \IfNoValueTF{#1}{%
    \propositionOLD
  }{%
    \propositionOLD
    \phantomsection
    \addcontentsline{toc}{subsubsection}{%
      \protect\numberline{\theproposition}\ #1}%
  }%
}{\endpropositionOLD}

% corollary
\let\corollaryOLD\corollary
\let\endcorollaryOLD\endcorollary
\RenewDocumentEnvironment{corollary}{o}{%
  \IfNoValueTF{#1}{%
    \corollaryOLD
  }{%
    \corollaryOLD
    \phantomsection
    \addcontentsline{toc}{subsubsection}{%
      \protect\numberline{\thecorollary}\ #1}%
  }%
}{\endcorollaryOLD}

% definition
\let\definitionOLD\definition
\let\enddefinitionOLD\enddefinition
\RenewDocumentEnvironment{definition}{o}{%
  \IfNoValueTF{#1}{%
    \definitionOLD
  }{%
    \definitionOLD
    \phantomsection
    \addcontentsline{toc}{subsubsection}{%
      \protect\numberline{\thedefinition}Definition\ #1}%
  }%
}{\enddefinitionOLD}
% ==================================================================== %
